\section{Results}

\subsection{Transformations and ephemeris composition}

All sampled interior comparisons pass their declared acceptance bounds
(Table~\ref{tab:validation}). Earth rotation and GMST agree with ERFA to below
$\RotationBound$ rad, and the maximum bias--precession--nutation matrix-element
residual is below $\BpnBound$. The terrestrial application's altitude
and azimuth agreement extends this check to the composition of time conversion,
sidereal rotation and local horizon geometry.

The time supplement additionally verifies civil leap-second handling and
interpolation of measured IERS data (Table~\ref{tab:time-extension}). The six UT1
comparisons pass a $1\,\mu$s bound. Rotations use full-JD inputs and agree
within $4\times10^{-9}$ rad with split-date references. One EOP coverage test
identifies a released boundary defect: a fractional day beyond the last sample
repeats that sample. In Tempoch 0.11.1 the last sample is MJD 60678,
whereas querying MJD 60678.0001 still returns its EOP values.
This outcome is classified as a failure rather than numerical agreement. Callers must reject nonfinite
epochs and enforce the closed interval between the first and last sample
before lookup. Stale DUT1 can affect UT1, Earth rotation and sidereal time;
stale polar-motion values can affect frame chains that consume them. This path is separate from the two
geometric applications.

\begin{table}[!tbp]
\centering
\caption{Supplemental time comparisons with split-date ERFA and IERS-derived
references. Maxima are absolute residuals; bounds are acceptance tolerances.
The separate EOP endpoint failure is not included among the numerical successes.}
\label{tab:time-extension}
\footnotesize
\begin{tabular}{@{}llrrr@{}}
\toprule
Observable & Unit & $n$ & Maximum & Bound \\
\midrule
UTC to TAI & s & 19 & 1.46e-11 & 1e-06 \\
UTC to TT & s & 19 & 2.29e-11 & 1e-06 \\
UTC to UT1 & s & 6 & 4.72e-08 & 1e-06 \\
IERS DUT1 & s & 18 & 0.00e+00 & 1e-12 \\
Polar motion $x_p$ & arcsec & 18 & 1.39e-17 & 1e-12 \\
Polar motion $y_p$ & arcsec & 18 & 0.00e+00 & 1e-12 \\
ERA & rad & 6 & 1.43e-09 & 4e-09 \\
GMST & rad & 6 & 1.43e-09 & 4e-09 \\
\bottomrule
\end{tabular}

\end{table}

The same-kernel ephemeris results combine kernel interpretation, body-center
chains, time handling and reference-frame conversion. Agreement over 1,002
interior epochs supports the implementation of this composed path. Exact
coverage endpoints return out-of-range errors, while points one second inside
succeed; this boundary behavior requires an interior margin for
these queries. Requests one day outside also return
coverage errors.

\begin{table*}[!tbp]
\centering
\caption{Independent comparisons for published Siderust v0.16.0. Time/frame
rows use ERFA; ephemeris rows use CSPICE with the same DE440s kernel and aligned
frames. Position residuals are Euclidean vector norms; BPN residuals are
maximum absolute matrix-element differences per epoch.
$^{*}$UTC rows compare full-JD \texttt{f64} outputs with a
$50\,\mu$s acceptance bound; zero denotes no resolved difference, not
exact time-scale agreement. Near J2000, one full-JD spacing is about
$40\,\mu$s.}
\label{tab:validation}
\footnotesize
\begin{tabular}{@{}llrrrrr@{}}
\toprule
Observable & Unit & $n$ & Median & p95 & p99 & Maximum \\
\midrule
Earth rotation angle & rad & 1000 & 9.492e-15 & 4.177e-14 & 5.373e-14 & 5.662e-14 \\
GMST (IAU 2006) & rad & 1000 & 9.492e-15 & 4.177e-14 & 5.373e-14 & 5.662e-14 \\
Mean obliquity & rad & 1000 & 5.551e-17 & 1.110e-16 & 1.665e-16 & 1.665e-16 \\
BPN matrix element & dimensionless & 1000 & 1.110e-16 & 2.241e-16 & 3.331e-16 & 3.993e-16 \\
UTC to TAI$^{*}$ & s & 10 & 0.000e+00 & 0.000e+00 & 0.000e+00 & 0.000e+00 \\
UTC to TT$^{*}$ & s & 10 & 0.000e+00 & 0.000e+00 & 0.000e+00 & 0.000e+00 \\
Case: azimuth & deg & 289 & 3.166e-11 & 8.235e-11 & 9.205e-11 & 9.325e-11 \\
Case: altitude & deg & 289 & 2.306e-11 & 2.866e-11 & 2.915e-11 & 2.929e-11 \\
Earth heliocentric & au & 1002 & 1.303e-16 & 2.570e-16 & 3.523e-16 & 5.561e-16 \\
Moon geocentric & km & 1002 & 4.757e-11 & 9.318e-11 & 1.304e-10 & 1.456e-10 \\
\bottomrule
\end{tabular}

\end{table*}

The analytical backend offers an alternative model path through the same
public ephemeris interface. Table~\ref{tab:models} reports its differences
from DE440s at the first 1,000 ephemeris epochs. These values characterize
the released analytical outputs relative to the SPK
backend over the sampled interval, including their frame-realization effects.
Separating them from same-kernel residuals makes the scientific
consequence of selecting an ephemeris backend explicit.

\begin{table*}[!tbp]
\centering
\caption{Backend-output differences from DE440s for v0.16.0 at 1,000 TT
epochs, eight days apart from J2000. Residuals are Euclidean vector norms and
include native analytical versus IAU 2006 frame-realization effects. Earth is
heliocentric and Moon geocentric; both paths compute geometric positions.}
\label{tab:models}
\footnotesize
\begin{tabular}{@{}llrrr@{}}
\toprule
Model & Unit & Median & p95 & Maximum \\
\midrule
Earth, VSOP87A & au & 3.423e-07 & 3.620e-07 & 3.725e-07 \\
Moon, ELP2000-82B & km & 3.218e-01 & 4.414e-01 & 4.717e-01 \\
\bottomrule
\end{tabular}

\end{table*}

\subsection{Generic surface observations}

The lunar and Mars comparisons pass their surface-geometry bounds
(Table~\ref{tab:candidate-surface}). The lunar rotation-matrix residual is
below $\LunarMatrixBound$ per element; distant and finite-target angles agree
within $\SurfaceAngleBound$ degrees. The finite-target comparison exercises
observer displacement and range together with orientation. A small rotation
residual can produce a larger Cartesian component difference at a large target
distance, explaining the different scales of the angle, range and component
rows.

All 80 additional analytic boundary cases pass, including 12 undefined-azimuth
cases returning the documented zero convention. The comparisons use separately
constructed reference geometry and the tolerances specified in the experimental
method. Small residuals are not estimates of physical surface or pointing
uncertainty.

The Earth-from-Moon result in Section~\ref{sec:lunar-example} demonstrates the
same pipeline with an ephemeris-derived target. Its independent agreement and
observable parallax connect the generic surface abstraction to a concrete
astronomical calculation. The cardinal diagnostic also verifies the released
Cartesian north/east/up convention.

\begin{table*}[!tbp]
\centering
\caption{Independent surface comparisons for published v0.16.0. Lunar rows
use analytical CSPICE \texttt{IAU\_MOON}; Mars rows check reference-ellipsoid
geometry. Matrix epochs and finite-target geometries are counted once each;
the ten finite-target geometries were each executed three times. Angular
residuals combine altitude and wrapped azimuth; component and basis rows use
maximum absolute component differences, and range uses an absolute difference.}
\label{tab:candidate-surface}
\footnotesize
\begin{tabular}{@{}llrrr@{}}
\toprule
Observable & Unit & $n$ & Maximum & Bound \\
\midrule
Lunar rotation matrix & dimensionless & 5 & 4.806e-11 & 1.0e-10 \\
Distant direction angle & deg & 30 & 4.499e-09 & 2.0e-08 \\
Finite target angle & deg & 10 & 3.827e-09 & 2.0e-08 \\
Finite target range & m & 10 & 7.510e-05 & 1.0e-03 \\
Finite Cartesian components & m & 10 & 1.756e-02 & 1.0e-01 \\
Mars surface position & m & 4 & 4.657e-10 & 1.0e-08 \\
Mars ENU basis & dimensionless & 4 & 1.110e-16 & 1.0e-14 \\
\bottomrule
\end{tabular}

\end{table*}

\subsection{Executable semantics and arithmetic cost}

The compiler rejects the five demonstrated dimension, frame, center, time-scale
and observer-body mismatches with the recorded \texttt{E0308} or
\texttt{E0277} diagnostics. Compatible operations execute successfully;
observer-site checks distinguish matching from mismatched runtime parameters.
These results connect the interface distinctions in Table~\ref{tab:guarantees}
to executable behavior.

Typed and scalar arithmetic costs are closely matched for the two measured
operations (Table~\ref{tab:cost}). The results support the practical use of
unit-bearing arithmetic in these kernels, with ordinary compiler optimization
and the same numerical outputs as their scalar counterparts.

\begin{table}[!tbp]
\centering
\caption{Paired arithmetic measurements for the published dependency graph:
30 rounds of $10^6$ operations per row. SD describes round-to-round variation
on the recorded host.}
\label{tab:cost}
\footnotesize
\begin{tabular}{@{}llrr@{}}
\toprule
Operation & Kernel & Median (ns/op) & SD (ns/op) \\
\midrule
km to m & scalar & 1.384 & 0.069 \\
km to m & typed & 1.384 & 0.055 \\
distance / time & scalar & 1.609 & 0.016 \\
distance / time & typed & 1.611 & 0.042 \\
\bottomrule
\end{tabular}

\end{table}

\begin{table}[!tbp]
\centering
\caption{Typed/scalar ratios within the original 30 paired rounds per operation.
The p5--p95 range is empirical round variability, not a confidence interval.}
\label{tab:ratios}
\footnotesize
\begin{tabular}{@{}lrrr@{}}
\toprule
Operation & Median ratio & p5 & p95 \\
\midrule
km to m & 1.0002 & 0.9595 & 1.0048 \\
distance / time & 1.0001 & 0.9936 & 1.0566 \\
\bottomrule
\end{tabular}

\end{table}

The paired medians in Table~\ref{tab:ratios} are close to unity. Variation across
rounds bounds the interpretation of differences at this scale on the recorded
host. These two kernels support no speed conclusion about full astronomical
workflows or other software systems.
